\section{Introduction}

The Transformer~\cite{vaswani2017attention} is the most-cited baseline in machine translation, and its WMT14 numbers --- 27.3 tokenized BLEU for Base on en-de, 38.1 for Big on en-fr --- are the reference point against which post-2017 progress is measured. Reproducing those numbers is harder than the paper's compactness suggests: their tokenized BLEU is not directly comparable to modern sacrebleu~\cite{post2018call}, the data preprocessing pipeline is under-specified, and the original training runs used 8-GPU configurations that no longer match the typical reproduction environment of a single accelerator. Most work that cites the Transformer paper quotes the numbers without re-running them. This paper re-runs them, on a single RTX 5090, and uses the resulting infrastructure to ask a sharper question.

\paragraph{The question.}
Modern empirical ML has converged on two largely-independent levers for improving a learning system: \emph{more capacity}~\cite{kaplan2020scaling, hoffmann2022chinchilla} and \emph{more, cleaner data}~\cite{penedo2023refinedweb, gururangan2020dont}. It is widely assumed that pulling either lever, holding the other roughly fixed, produces monotone gains. Our study contradicts that assumption in a sharp form. On WMT14 en-fr, scaling Base (60M) to Big (209M) \textbf{changes sign as a function of data quality}: on a strictly-filtered 9.3M corpus, Big beats Base by $+0.56$ BLEU; on a loosely-filtered 30M corpus drawn from the same source pool, Big falls below Base by $-0.87$ BLEU. The same architectural intervention pays in one regime and costs in another; the regime is determined entirely by data quality.

\paragraph{What we do.}
We construct a clean two-by-two design: data quality (strict-filter v1, loose-filter v2) $\times$ model capacity (Base, Big), and report all four trained-from-scratch results on the same single-GPU pipeline. We extend the design with cross-language replication (WMT14 en-de, where the same Base $\geq$ Big ordering holds at 4.5M scale) and a post-training intervention (SFT on the top-1M of the noisy 30M, scored by CometKiwi-22~\cite{rei2022cometkiwi}). The SFT result is a clean negative finding that tightens the central claim: data quality is not a single number --- it decomposes into translation correctness \emph{and} alignment to the eval distribution, and naive top-K-by-QE filtering optimizes only the first.

\paragraph{Contributions.}
\begin{enumerate}
\item A 2$\times$2 ablation establishing that capacity return is sign-conditioned on data quality ($+0.56$ BLEU on clean data, $-0.87$ on noisy data, same en-fr corpus pool, same SPM, same pipeline; replicated qualitatively on en-de).
\item A post-training negative result showing that quality-filtered SFT, at scale, degrades newstest BLEU by 1.5--2.3, with the larger model degrading more. The failure is mechanistically distinct from the pretraining-on-noise failure of \S\ref{sec:enfr_v2} but converges to the same 33--34 BLEU ceiling, suggesting that incorrect supervision of \emph{any} form bounds the achievable test accuracy by approximately the same amount.
\item A reproducible single-GPU Transformer pipeline that hits 35.31 / 35.87 sacrebleu on WMT14 en-fr Base/Big and 24.04 on en-de Base, with full open-source code, configurations, four published HuggingFace checkpoints, and persistent training logs. The whole study reproduces in approximately 30 GPU-hours on a single 5090.
\item A methodology note on loss-spike-guard trigger rates as a diagnostic, with evidence that they cannot be interpreted as a pure data-noise metric: the trigger rate of a fixed threshold is a joint function of data noise and the model's own convergence depth, so cross-capacity comparisons (Big on clean vs.\ Big on noisy) are confounded.
\end{enumerate}

\paragraph{Scope and limitations.}
We run single seeds per configuration. The variance of any one BLEU number we report is unmeasured, but the $+0.56$ / $-0.87$ / $-1.55$ / $-2.33$ effects we make claims on are large relative to the typical Transformer-MT seed-variance band of ${\sim}0.2$--$0.4$ BLEU, and the cross-language replication argues for the directional findings. We use sacrebleu 13a throughout; we do not run human evaluation. We test only WMT14 en-fr and en-de; generalization to morphologically richer or non-Latin pairs (e.g., en-cs, en-ru, en-zh) is not in scope.

\paragraph{Roadmap.}
\S\ref{sec:related} reviews related work. \S\ref{sec:setup} describes the single-GPU setup. \S\ref{sec:enfr_v1}--\S\ref{sec:enfr_v2} give the en-fr v1 success and v2 failure runs that constitute the data-quality axis. \S\ref{sec:ende} generalizes the pattern to en-de. \S\ref{sec:sft} reports the quality-filtered-SFT negative result and the multi-dimensional view of ``data quality'' it implies. \S\ref{sec:methodology} is the methodology note on spike-guard interpretation. \S\ref{sec:conclusion} concludes.
